\scanpage{143}

montre bien que dans certains cas, la validité de l'opération $A_0\longmapsto A_0^{\unsure{\nu|v|\vee}}$ déborde de celle de l'axiomatique $\mathrm{To}\,1:\mathrm{To}\,5$ développée jusqu'à présent.

Le cas suivant serait
\[
 n=7,\qquad Z'_n\simeq\mathbb Z/3\mathbb Z
 \qquad \deleted{\unsure{\simeq\mathbb Z|=}}\quad
 ((\mathbb Z/7\mathbb Z)^\times\simeq\mathbb Z/6\mathbb Z)
\]
engendré par n'importe quel
\[
 \unsure{\zeta|\nu|\xi}\in\mathbb F_7^\times\setminus\deleted{\unsure{\{\pm\}|\{1\}}}\{1,-1\}
 =\mathbb F_7^{\deleted{\unsure{\times|*}}}\setminus\{0,1,-1\}
\]
p.ex. $\nu=2\bmod 7$, dont l'inverse est $\nu'=4\bmod 7$,

est-il vrai que pour tout 2-polyèdre \deleted{\unsure{rég|rég.}} triangulé \inserted{$(S,A_0,F_0)$} dont les sommets sont \inserted{tous} d'ordre 7, \unsure{l'|le} \deleted{\unsure{g|gr}} gr. pour $(S,A_0^{(2)}\subset\mathfrak P_2(S))$ correspond encore à un tel 2-polyèdre ? Même question pour $n=8$, avec $Z'_8\simeq\mathbb Z/2\mathbb Z$, ou \inserted{$n=10$ avec} $Z'_{10}\simeq\mathbb Z/2\mathbb Z$ (pour $n=9$ on a $Z'_9\simeq\mathbb Z/4\mathbb Z$ -- en fait pour
\[
\left\{\begin{aligned}
&n\notin\{\underbrace{2,3,4,6}_{Z'_n=1},
 \underbrace{5,8,10\unsure{\}|{}},12}_{Z'_n\simeq\mathbb Z/2\mathbb Z})
 \quad\text{i.e.}\quad n\geq18\ \text{ou}\ n\in\{
 \underbrace{7}_{Z'_n\simeq\mathbb Z/3\mathbb Z},
 \underbrace{9}_{Z'_n\in\mathbb Z/4\mathbb Z},11\},\ \text{on a}\\
&Z'_n\ \text{d'ordre}\ \geq3\ \text{-- ordre 3 pour}\ n=7\ \text{et }14,\ \text{ordre }4\\
&\text{pour }n=9\ \text{et }18\ \inserted{\text{et pour }n=16,20}\ \text{-- ordre 5 pour }n=11\ \text{et }22,\ \text{-- ordre 6 pour}\\
&n=13\ \text{et }26\ \text{-- }\inserted{\text{(ordre 12 pour }n=15\ \text{ou }30\text{)}}\
  \deleted{\text{\unsure{ordre|ordre 8}}}\ \text{ordre}\\
&8\ \text{pour }n=17\ \text{ou }34,\ \deleted{\text{\unsure{10 pour|ordre}}}\
 \inserted{\text{9 pour}}\ 19\ \text{et }38,
\end{aligned}\right.
\]

Mais poursuivons sur la lancée \unsure{des|de} l'axiomatique des $\mathrm{To}\,1\ldots$, on trouve
\[
 \text{(To 7)}\quad \text{L'homomorphisme}\quad Z\hookrightarrow\Lambda'_n\quad\text{est surjectif.}\tag*{(38)}
\]
i.e. un isomorphisme.

% archive sheet number: 142
